%! Characteristics of Logarithmic Functions — Unit 8, Lesson 8.0 — Guided Notes
\documentclass[10pt]{article}
\usepackage{algebra2-article}
\usepackage{algebra2-boxes}
\newcommand{\TallMath}[1]{$\displaystyle #1\rule[-1.4em]{0pt}{3.2em}$}
\newcommand{\lgrid}[4]{%
  \draw[step=1, linegray!40, very thin] (#1,#3) grid (#2,#4);
  \draw[->, semithick, charcoal] (#1,0)--(#2,0) node[below right, font=\scriptsize]{$x$};
  \draw[->, semithick, charcoal] (0,#3)--(0,#4) node[left, font=\scriptsize]{$y$};}
% pgfmath has no log_b, so every logarithmic curve is drawn by parameterizing on y:
%   y = log_b(x-h) + k   <=>   x = b^(y-k) + h.   Args: {b}{h}{k}{ymin}{ymax}
\newcommand{\logcurve}[5]{\draw[very thick, forest, domain=#4:#5, samples=120, smooth]
  plot ({pow(#1,\x-(#3))+(#2)},{\x});}

\begin{document}

\pageheader{Unit 8, Lesson 8.0}{Guided Notes}
\namedateperiod

\vspace{0.08in}

\begin{objectivebox}
\textbf{By the end of this lesson, I will be able to\dots}
\begin{itemize}\setlength{\itemsep}{2pt}
  \item find a logarithmic function's \emph{domain} by setting its \emph{argument} $>0$, and name the \emph{vertical asymptote} that sits at the edge of that domain;
  \item read the \emph{domain}, \emph{range}, \emph{asymptote}, \emph{intercepts}, \emph{increasing/decreasing} and \emph{positive/negative} intervals, \emph{extrema}, and \emph{end behavior} off a logarithmic graph with base $b>1$ or $0<b<1$;
  \item explain why $y=\log_b x$ is the \emph{inverse} of $y=b^x$ --- so their domains and ranges \emph{swap}, and the exponential's horizontal asymptote becomes the logarithm's vertical one.
\end{itemize}
\end{objectivebox}

\vspace{0.05in}

\begin{vocabbox}
\small
Fill in each term as we build it together.
\par\vspace{2pt}   % \par is required: \termblanklong's \noindent is a no-op mid-paragraph
\termblanklong{Logarithmic function}
\termblanklong{Base $b$}
\termblanklong{Argument}
\termblanklong{Domain restriction (from the argument)}
\termblanklong{Vertical asymptote}
\termblanklong{Inverse relationship (exponential $\leftrightarrow$ logarithmic)}
\end{vocabbox}

\begin{hookbox}
Here is the graph Unit~7 ended with, $y=2^x$, and beside it the brand-new graph this unit is built
on, $y=\log_2 x$.
\begin{center}
\begin{tikzpicture}[scale=0.38]
  % y = 2^x
  \begin{scope}
    \lgrid{-4}{5}{-4}{5}
    \draw[navy, dashed, semithick] (-4,0)--(5,0);
    \begin{scope}
      \clip (-4,-4) rectangle (5,5);
      \draw[very thick, forest, domain=-4:2.32, samples=90, smooth] plot (\x,{pow(2,\x)});
    \end{scope}
    \fill[charcoal] (0,1) circle (3pt);
    \node[charcoal, font=\scriptsize] at (0.5,-5.4) {$y=2^x$};
  \end{scope}
  % y = log_2 x
  \begin{scope}[xshift=13cm]
    \lgrid{-4}{5}{-4}{5}
    \draw[navy, dashed, semithick] (0,-4)--(0,5);
    \begin{scope}
      \clip (-4,-4) rectangle (5,5);
      \logcurve{2}{0}{0}{-4}{2.32}
    \end{scope}
    \fill[charcoal] (1,0) circle (3pt);
    \node[charcoal, font=\scriptsize] at (0.5,-5.4) {$y=\log_2 x$};
  \end{scope}
\end{tikzpicture}
\end{center}
\begin{itemize}\setlength{\itemsep}{1pt}
  \item The second graph is the first one \textbf{tipped over}. Which line would you have to fold the page along to turn one into the other? \writeline
  \item The exponential graph never touched the $x$-axis. The new graph never touches the $y$-axis. What happened to that flat line when the page was folded? \writeline
\end{itemize}
Unit~7 asked, ``I have an exponent --- what is the value?'' This unit asks it the other way:
\textbf{``I have the value --- what was the exponent?''} Today's job is to learn to read the graph
that answers that question. \emph{Naming} it comes first; \emph{computing} with it is Lesson~8.1.
\end{hookbox}

\begin{notesbox}{1. What a Logarithmic Function Is --- and Where It Is Allowed to Exist}
Warm-Up items (d) and (e) already asked the backwards question: \emph{$2$ to what power gives $32$?}
A \textbf{logarithm} is just a name for that answer.
\[
  \log_b x = \text{the exponent you put on } b \text{ to get } x .
\]
Two parts have names, and both matter today:
\begin{itemize}[leftmargin=1.5em, itemsep=3pt]
  \item The \textbf{base} $b$ is the small number written low --- the number being raised to a power.
        Just as in Unit~7, it must satisfy $b>0$ and $b\ne\blank{0.6cm}$.
  \item The \textbf{argument} is the expression \blank{4.6cm} the logarithm --- the value you are
        trying to \emph{reach}.
\end{itemize}
Read these straight off the definition (no calculator):
\begin{center}
\renewcommand{\arraystretch}{1.4}
\begin{tabular}{@{}l l l@{}}
$\log_2 8=$ \blank{1.2cm} because $2^{\square}=8$ &
$\log_2 1=$ \blank{1.2cm} &
$\log_3 9=$ \blank{1.2cm} \\
$\log_2 \tfrac14=$ \blank{1.2cm} &
$\log_5 5=$ \blank{1.2cm} &
$\log_2 (-8)=$ \blank{2.6cm} \\
\end{tabular}
\end{center}

\vspace{3pt}
\textbf{The rule that runs this unit.} That last box is the important one. From Warm-Up item~1(f):
no exponent makes $2^{\square}$ negative or zero, so there is nothing for $\log_2(-8)$ to name. It is
\blank{2.6cm}.
\begin{center}
\fbox{\parbox{0.92\linewidth}{\centering \textbf{The argument of a logarithm must be strictly
positive.} \\ For $y=\log_b(\text{argument})$, set the \textbf{argument} $>0$ and solve.}}
\end{center}
\begin{itemize}[leftmargin=1.5em, itemsep=3pt]
  \item $g(x)=\log_2(x-1)$: \; solve $x-1>0$, so the domain is \blank{3.4cm}.
  \item $h(x)=\log_3(x+4)$: \; solve $x+4>0$, so the domain is \blank{3.4cm}.
  \item $k(x)=\log(5-x)$: \; solve $5-x>0$, so the domain is \blank{3.4cm}. \\
        \emph{Notice:} this one runs to the \textbf{left}. A negative in front of the $x$ flips which
        way the graph goes --- exactly as it did with radicals in Unit~6.
\end{itemize}

\vspace{2pt}
\textbf{One character different --- and it matters.} A square root allowed radicand $\ge 0$, so the
graph got to \emph{stand on} its edge point. A logarithm demands argument $>\blank{0.6cm}$
\textbf{strictly}, so the edge is \emph{missing} --- the graph runs alongside it forever without ever
arriving. That missing edge is a \blank{4.2cm}, the Unit~5 idea returning in a brand-new family.
\end{notesbox}

\begin{notesbox}{2. The Logarithmic Parent, $b>1$: $y=\log_2 x$}
\begin{center}
\begin{minipage}[c]{0.40\linewidth}
\centering
\renewcommand{\arraystretch}{1.25}
\begin{tabular}{|c|c|}
\hline
$x$ & $y=\log_2 x$ \\ \hline
$\tfrac14$ & \blank{1.0cm} \\ \hline
$\tfrac12$ & \blank{1.0cm} \\ \hline
$1$ & \blank{1.0cm} \\ \hline
$2$ & \blank{1.0cm} \\ \hline
$4$ & \blank{1.0cm} \\ \hline
$8$ & \blank{1.0cm} \\ \hline
\end{tabular}
\end{minipage}
\begin{minipage}[c]{0.55\linewidth}
\centering
\begin{tikzpicture}[scale=0.46]
  \lgrid{-2}{9}{-4}{4}
  \draw[navy, dashed, semithick] (0,-4)--(0,4);
  \begin{scope}
    \clip (-2,-4) rectangle (9,4);
    \logcurve{2}{0}{0}{-4}{3.17}
  \end{scope}
  \fill[charcoal] (1,0) circle (2.6pt);
  \fill[charcoal] (2,1) circle (2.6pt);
  \fill[charcoal] (4,2) circle (2.6pt);
  \fill[charcoal] (8,3) circle (2.6pt);
  \node[navy, font=\scriptsize, right] at (0.15,-3.4) {$x=0$};
\end{tikzpicture}
\end{minipage}
\end{center}

\vspace{2pt}
The graph hugs the $y$-axis on the way down and climbs forever on the right --- but \emph{slowly}.
\begin{enumerate}[leftmargin=1.7em, itemsep=3pt]
  \item \textbf{Domain:} \blank{3.2cm} \quad \textbf{Range:} \blank{3.2cm}
  \item \textbf{Vertical asymptote:} the line \blank{1.6cm}. The graph gets arbitrarily close but
        \emph{never} touches it, because $x=0$ is \blank{3.6cm} the domain.
  \item \textbf{Intercepts:} $x$-intercept \blank{1.6cm} --- every parent log graph passes through
        this point, because $b^0=\blank{0.6cm}$ for \emph{any} base. \\
        $y$-intercept: \blank{2.0cm}, because $x=0$ is not in the domain.
  \item \textbf{Increasing / decreasing:} increasing on \blank{3.0cm}; it is \emph{never}
        \blank{2.4cm}.
  \item \textbf{Positive / negative:} $y<0$ on \blank{2.4cm}; \quad $y>0$ on \blank{2.4cm}.
  \item \textbf{Extrema:} \blank{2.0cm}. The graph falls forever near the asymptote and climbs
        forever to the right, so there is no absolute maximum and no absolute minimum.
  \item \textbf{End behavior:} as $x\to+\infty$, $y\to\blank{1.2cm}$ --- but \emph{very} slowly
        ($x$ must \emph{double} just to push $y$ up by $1$). \\
        And on the other side, the input cannot run off to $-\infty$; it runs down to the asymptote
        instead: as $x\to 0^{+}$, $y\to\blank{1.2cm}$.
\end{enumerate}
\end{notesbox}

\begin{notesbox}{3. The Other Case, $0<b<1$: $y=\log_{1/2} x$}
\begin{center}
\begin{minipage}[c]{0.40\linewidth}
\centering
\renewcommand{\arraystretch}{1.25}
\begin{tabular}{|c|c|}
\hline
$x$ & $y=\log_{1/2} x$ \\ \hline
$\tfrac14$ & \blank{1.0cm} \\ \hline
$1$ & \blank{1.0cm} \\ \hline
$2$ & \blank{1.0cm} \\ \hline
$4$ & \blank{1.0cm} \\ \hline
$8$ & \blank{1.0cm} \\ \hline
\end{tabular}
\end{minipage}
\begin{minipage}[c]{0.55\linewidth}
\centering
\begin{tikzpicture}[scale=0.46]
  \lgrid{-2}{9}{-4}{4}
  \draw[navy, dashed, semithick] (0,-4)--(0,4);
  \begin{scope}
    \clip (-2,-4) rectangle (9,4);
    \logcurve{0.5}{0}{0}{-3.17}{4}
  \end{scope}
  \fill[charcoal] (1,0) circle (2.6pt);
  \fill[charcoal] (2,-1) circle (2.6pt);
  \fill[charcoal] (4,-2) circle (2.6pt);
  \fill[charcoal] (8,-3) circle (2.6pt);
\end{tikzpicture}
\end{minipage}
\end{center}

\vspace{2pt}
A base between $0$ and $1$ is the mirror of a base above $1$ --- the same way \emph{decay} mirrored
\emph{growth} in Unit~7. Fill in only what \textbf{changed}, and notice how much did not.
\begin{center}
\renewcommand{\arraystretch}{1.3}
\begin{tabularx}{\linewidth}{@{}l X X@{}}
\hline
 & \textbf{$y=\log_2 x$} \; $(b>1)$ & \textbf{$y=\log_{1/2} x$} \; $(0<b<1)$ \\
\hline
Domain & $(0,\infty)$ & \blank{2.4cm} \\
Range & all reals & \blank{2.4cm} \\
Vertical asymptote & $x=0$ & \blank{2.4cm} \\
$x$-intercept & $(1,0)$ & \blank{2.4cm} \\
$y$-intercept & none & \blank{2.4cm} \\
Increasing or decreasing & increasing & \blank{2.4cm} \\
As $x\to+\infty$ & $y\to+\infty$ & \blank{2.4cm} \\
As $x\to 0^{+}$ & $y\to-\infty$ & \blank{2.4cm} \\
\hline
\end{tabularx}
\end{center}
Four rows are \emph{identical} in both columns. Which four --- and what single fact do all four of
them come from? \writeline
\end{notesbox}

\begin{notesbox}{4. Where This Graph Came From --- the Inverse of the Exponential}
Warm-Up item~3 already showed it: the table for $y=2^x$ and the table for $y=\log_2 x$ are the
\emph{same table with the rows traded}. Raising $2$ to a power and asking ``$2$ to what power?''
\textbf{undo each other} --- they are \textbf{inverses} --- and on a graph, undoing looks like a
\textbf{reflection over the line $y=x$}.
\begin{center}
\begin{tikzpicture}[scale=0.44]
  \lgrid{-4}{6}{-4}{6}
  \draw[navy, dashed, semithick] (-4,-4)--(6,6);
  \node[navy, font=\scriptsize, above left] at (5.6,5.6) {$y=x$};
  \draw[goldacc, dashed, semithick] (-4,0)--(6,0);
  \draw[forest, dashed, semithick] (0,-4)--(0,6);
  \begin{scope}
    \clip (-4,-4) rectangle (6,6);
    \draw[very thick, goldacc, domain=-4:2.58, samples=90, smooth] plot (\x,{pow(2,\x)});
    \logcurve{2}{0}{0}{-4}{2.58}
  \end{scope}
  \fill[charcoal] (0,1) circle (2.6pt) node[above left, font=\scriptsize]{$(0,1)$};
  \fill[charcoal] (1,0) circle (2.6pt) node[below right, font=\scriptsize]{$(1,0)$};
  \node[goldacc, font=\scriptsize, right] at (1.9,5.2) {$y=2^x$};
  \node[forest, font=\scriptsize, right] at (4.2,2.4) {$y=\log_2 x$};
\end{tikzpicture}
\end{center}
\begin{itemize}[leftmargin=1.5em, itemsep=3pt]
  \item Reflecting over $y=x$ \textbf{swaps the coordinates} of every point: the point $(8,3)$ on
        $y=\log_2 x$ came from $(\blank{0.6cm},\ \blank{0.6cm})$ on $y=2^x$. So the two functions
        also swap their \blank{2.0cm} and their \blank{2.0cm}.
  \item That is the \emph{deeper} reason for today's domain rule. The domain of $y=\log_2 x$ is
        $(0,\infty)$ because that is the \blank{2.0cm} of $y=2^x$ --- an exponential never produces a
        negative output, so a logarithm never accepts a negative input.
  \item Read it the other way too: the \emph{range} of $y=\log_2 x$ is all reals because that was the
        \blank{2.0cm} of $y=2^x$. Any real number is a legal exponent.
  \item \textbf{The asymptote turns a corner.} The exponential's \emph{horizontal} asymptote was the
        line $y=0$. Reflect that line over $y=x$ and you get the line \blank{1.6cm} --- the
        logarithm's \emph{vertical} asymptote. It is the same boundary, seen sideways.
  \item \textbf{Same for the intercepts.} $y=2^x$ had a $y$-intercept at $(0,1)$ and \emph{no}
        $x$-intercept, ever. Trade the coordinates: $y=\log_2 x$ has an $x$-intercept at
        \blank{1.4cm} and \emph{no} \blank{2.0cm}, ever. One fact, reflected.
  \item \textbf{Why is no restriction needed here?} In Unit~6, $y=x^2$ had to be cut down to $x\ge0$
        before $\sqrt{x}$ could undo it, because a parabola sends two inputs to the same output. An
        exponential is always increasing (or always decreasing), so it \emph{never} repeats an
        output. All of it can be reflected --- \blank{4.4cm} needed.
\end{itemize}
\emph{Hold on to this.} It is the reason logarithmic graphs look the way they do, and Lessons~8.1
and~8.4 turn it into a method for solving equations.
\end{notesbox}

\begin{notesbox}{5. The Full Feature List --- and How the Two Families Differ}
Now read \emph{every} characteristic off the anchor $f(x)=\log_2(x-1)$.
\begin{center}
\begin{tikzpicture}[scale=0.50]
  \lgrid{-2}{10}{-4}{4}
  \draw[navy, dashed, semithick] (1,-4)--(1,4);
  \begin{scope}
    \clip (-2,-4) rectangle (10,4);
    \logcurve{2}{1}{0}{-4}{3.17}
  \end{scope}
  \fill[charcoal] (2,0) circle (2.6pt) node[below right, font=\scriptsize]{$(2,0)$};
  \fill[charcoal] (3,1) circle (2.6pt) node[above left, font=\scriptsize]{$(3,1)$};
  \fill[charcoal] (5,2) circle (2.6pt) node[above left, font=\scriptsize]{$(5,2)$};
  \fill[charcoal] (9,3) circle (2.6pt) node[above left, font=\scriptsize]{$(9,3)$};
  \node[navy, font=\scriptsize, right] at (1.15,-3.4) {$x=1$};
\end{tikzpicture}
\end{center}
\begin{enumerate}[leftmargin=1.7em, itemsep=3pt]
  \item \textbf{Domain} (set the argument $x-1>0$): \blank{3.4cm} \quad \textbf{Range:}
        \blank{3.4cm}
  \item \textbf{Vertical asymptote:} \blank{1.8cm} --- and notice the domain \emph{starts} at that
        same number, with a \textbf{round} bracket.
  \item \textbf{$x$-intercept / zero:} \blank{1.8cm} \quad \textbf{$y$-intercept:} \blank{1.8cm}
  \item \textbf{Increasing / decreasing:} increasing on \blank{3.0cm} --- the entire domain.
  \item \textbf{Positive / negative:} $f(x)<0$ on \blank{2.6cm}; \quad $f(x)>0$ on \blank{2.6cm}.
        \emph{Careful:} neither interval may include inputs outside the domain.
  \item \textbf{Extrema:} \blank{2.6cm} --- no absolute maximum and no absolute minimum.
  \item \textbf{End behavior:} as $x\to+\infty$, $f(x)\to\blank{1.2cm}$. \; As $x\to 1^{+}$,
        $f(x)\to\blank{1.2cm}$.
\end{enumerate}

\vspace{4pt}
\textbf{Compare and contrast} the logarithmic family with the exponential family it came from:
\begin{center}
\renewcommand{\arraystretch}{1.3}
\begin{tabularx}{\linewidth}{@{}l X X@{}}
\hline
 & \textbf{Exponential} $y=b^x$ & \textbf{Logarithmic} $y=\log_b x$ \\
\hline
Domain & all reals & \blank{2.4cm} \\
Range & \blank{2.4cm} & \blank{2.4cm} \\
Asymptote & horizontal, \blank{1.6cm} & vertical, \blank{1.6cm} \\
$x$-intercept & \blank{1.6cm} & \blank{1.6cm} \\
$y$-intercept & $(0,1)$ & \blank{1.6cm} \\
Extrema & none & \blank{1.6cm} \\
Inverse of & $y=\log_b x$ & \blank{1.8cm} \\
\hline
\end{tabularx}
\end{center}
And against the families you already know: a \textbf{polynomial} was defined everywhere; a
\textbf{rational} function lost scattered inputs and gained asymptotes; a \textbf{radical} function
lost half the number line but got to \emph{stand on} its edge. A \textbf{logarithmic} function loses
half the number line \emph{and} the edge itself --- the boundary is an asymptote, so the graph
approaches it forever and never arrives.
\end{notesbox}

\begin{practicebox}
The graph below is $p(x) = \log_2(4-x)$. Read each characteristic and write it on the line.
\begin{center}
\begin{tikzpicture}[scale=0.48]
  \lgrid{-6}{7}{-3}{4}
  \draw[navy, dashed, semithick] (4,-3)--(4,4);
  \begin{scope}
    \clip (-6,-3) rectangle (7,4);
    \draw[very thick, forest, domain=-3:3.32, samples=120, smooth] plot ({4-pow(2,\x)},{\x});
  \end{scope}
  \fill[charcoal] (3,0) circle (2.6pt) node[below left, font=\scriptsize]{$(3,0)$};
  \fill[charcoal] (2,1) circle (2.6pt);
  \fill[charcoal] (0,2) circle (2.6pt) node[above right, font=\scriptsize]{$(0,2)$};
  \fill[charcoal] (-4,3) circle (2.6pt);
\end{tikzpicture}
\end{center}
\begin{enumerate}[leftmargin=1.7em, itemsep=3pt]
  \item Set the argument $>0$: $4-x>0$, so the \textbf{domain} is \blank{3.4cm}. \quad
        \textbf{Range:} \blank{3.0cm}
  \item \textbf{Vertical asymptote:} \blank{1.6cm}. \quad Is $p$ increasing or decreasing on its
        domain? \blank{2.4cm}
  \item \textbf{$x$-intercept:} \blank{1.6cm} \quad \textbf{$y$-intercept:} \blank{1.6cm}
  \item This graph runs to the \textbf{left} while the parent $y=\log_2 x$ runs to the right. Which
        part of the equation caused that? \blank{5.0cm}
  \item This graph \emph{has} a $y$-intercept, but the parent $y=\log_2 x$ has none. What changed?
        \blank{5.0cm}
\end{enumerate}
\end{practicebox}

\end{document}
